\section{Matter: Spin from Topology, and a Chiral Fermion (P4)}
\label{sec:matter}

Vibe Theory builds matter from the same tones as everything else, with no new primitive. A particle is a persistent pattern of the mesh, and its spin is meant to come from the shape of the connections, not from an added ingredient. P4 tests this for the fermion.

A word on what a spinor is here, since the base tone is strictly ternary. A spinor is not a richer tone on a single vibe. It is a graded complex of ordinary ternary vibes: the vertices of a small patch are base vibes, the edges are relational vibes (the notes), and the faces are higher relational vibes, with the spinor being the tuple of their ternary tones knit together by the discrete exterior derivative. Every component is a vibe, the component count is fixed by the cell combinatorics rather than chosen by hand, and the gamma matrices are the note operations that raise or lower relational order. The complex phase a spinor uses is emergent, not primitive, carried by a finite clock-tone or the parity of the beat. So nothing richer than ternary is added: the spinor is a pattern of ternary vibes, and "a richer tone" would be the wrong way to say it.

A caveat on what the testbed measures. The results below are computed at the effective level, on a cell complex where the spinor is already a Dirac structure, with cell values represented as real or complex numbers (the emergent representation). This validates the spinor's behaviour, its spin and its chirality, at that level. It does not yet build the spinor bottom-up from an explicit cluster of ternary vibes evolving under the rule, which is the open construction for the matter sector.

\subsection{The monist spinor and spin from topology}

We build the Kahler-Dirac operator $D = d + \delta$ on the cochain spaces of a triangulated mesh. Its zero modes are the harmonic forms, whose count is a topological invariant, the sum of the Betti numbers. We built surfaces of known topology and counted.

\begin{result}[Spin from topology]
The Kahler-Dirac zero-mode count equals the surface Betti sum exactly: a disk gives $1$, a cylinder gives $2$, a torus gives $4$. Every component of each zero mode is a cell tone, so the spinor is built entirely from the substrate. The statement that a tone is a small spinor is realized, and the spinor's zero modes are a topological invariant of the mesh, which is the spin-from-topology reading (in the spirit of Bombelli et al., 1987).
\end{result}

\subsection{The chirality wall, threaded}

The hardest matter fact is that a single chiral (handed) fermion cannot live on a naive lattice. The Nielsen-Ninomiya theorem forces doublers: a local, hermitian, chirally symmetric lattice Dirac operator comes with equal numbers of left- and right-handed copies. We measured this and its resolutions in 2D momentum space.

\begin{result}[Overlap threads Nielsen-Ninomiya]
The naive lattice operator has $4$ species (the $2^2$ doublers). The Wilson operator removes them to $1$ but breaks chiral symmetry (a large Ginsparg-Wilson residual). The overlap operator $D = 1 + \gamma_5\, \mathrm{sign}(H_W)$ has $1$ species and an exact lattice chiral symmetry, with Ginsparg-Wilson residual $6 \times 10^{-16}$, machine zero. A single chiral fermion exists on the mesh, with exact chiral symmetry, threading the no-go theorem.
\end{result}

So matter is monist and its spin is topological, and the deepest obstruction to chiral matter is passed at the kinematic level. The remaining wall, a fully interacting Weyl-projected chiral gauge theory, is open in physics itself and is discussed with the forces.
